\documentclass{article}

\usepackage{neurips_2024}

\usepackage[utf8]{inputenc}
\usepackage[T1]{fontenc}
\usepackage{hyperref}
\usepackage{url}
\usepackage{booktabs}
\usepackage{amsfonts}
\usepackage{amsmath}
\usepackage{nicefrac}
\usepackage{microtype}
\usepackage{xcolor}
\usepackage{graphicx}
\usepackage{longtable}

\title{VAE-Based Unsupervised Clustering for Hybrid Language Music Tracks: A Comprehensive Multi-Modal Analysis}

\author{%
  Student Name \\
  Department of Computer Science and Engineering\\
  BRAC University\\
  Dhaka, Bangladesh \\
  \texttt{student.id@g.bracu.ac.bd}
}

\begin{document}

\maketitle

\begin{abstract}
This paper presents a comprehensive unsupervised learning pipeline for clustering music tracks using Variational Autoencoders (VAE). We work with a hybrid language music dataset containing both English and Bangla songs, extracting latent representations from audio features (log-mel spectrograms, MFCCs) and lyrics embeddings (TF-IDF). We implement and compare three levels of complexity aligned with progressive task requirements: (1) \textbf{Easy Task}: Basic VAE with K-Means clustering, t-SNE/PCA visualization, and baseline comparison; (2) \textbf{Medium Task}: Convolutional VAE for spectrograms, hybrid audio+lyrics feature representation, and multiple clustering algorithms (K-Means, Agglomerative, DBSCAN, GMM); and (3) \textbf{Hard Task}: Beta-VAE and Conditional VAE (CVAE) for disentangled representations, multi-modal fusion combining audio, lyrics, and genre information, with detailed cluster purity analysis. Our experiments on 3,213 audio clips spanning 82 genres demonstrate that Beta-VAE achieves the highest Silhouette Score of 0.366, while ConvVAE achieves the best supervised metrics (NMI=0.209, Purity=0.380). We provide comprehensive evaluations using Silhouette Score, Calinski-Harabasz Index, Davies-Bouldin Index, Adjusted Rand Index (ARI), Normalized Mutual Information (NMI), and Cluster Purity, along with detailed visualizations including t-SNE plots, cluster distribution over genres, and genre-cluster heatmaps.
\end{abstract}

\section{Introduction}

Music information retrieval (MIR) and automatic music categorization have become increasingly important with the exponential growth of digital music libraries. Traditional genre classification relies on supervised learning with manually labeled data, which is expensive and subjective. Unsupervised clustering offers an alternative approach that can discover natural groupings in music data without explicit labels.

Variational Autoencoders (VAEs) have emerged as powerful generative models capable of learning meaningful latent representations. Unlike traditional autoencoders, VAEs impose a probabilistic structure on the latent space, encouraging smooth and interpretable representations that are well-suited for downstream clustering tasks.

\subsection{Task Overview}

This project implements a three-tiered approach with increasing complexity:

\textbf{Easy Task:}
\begin{itemize}
    \item Implement a basic VAE for feature extraction from music data
    \item Use a small hybrid language music dataset (English + Bangla songs)
    \item Perform clustering using K-Means on latent features
    \item Visualize clusters using t-SNE and PCA
    \item Compare with baseline (PCA + K-Means) using Silhouette Score and Calinski-Harabasz Index
\end{itemize}

\textbf{Medium Task:}
\begin{itemize}
    \item Enhance VAE with convolutional architecture for spectrograms/MFCC features
    \item Include hybrid feature representation: audio + lyrics embeddings
    \item Experiment with multiple clustering algorithms: K-Means, Agglomerative Clustering, DBSCAN
    \item Evaluate clustering quality using Silhouette Score, Davies-Bouldin Index, and Adjusted Rand Index (with partial labels)
    \item Compare results across methods and analyze why VAE representations perform better/worse than baselines
\end{itemize}

\textbf{Hard Task:}
\begin{itemize}
    \item Implement Conditional VAE (CVAE) or Beta-VAE for disentangled latent representations
    \item Perform multi-modal clustering combining audio, lyrics, and genre information
    \item Quantitatively evaluate using Silhouette Score, Normalized Mutual Information (NMI), Adjusted Rand Index (ARI), and Cluster Purity
    \item Provide detailed visualizations: latent space plots, cluster distribution over languages/genres, reconstruction examples
    \item Compare VAE-based clustering with PCA + K-Means, Autoencoder + K-Means, and direct spectral feature clustering
\end{itemize}

\subsection{Contributions}

In this work, we implement a comprehensive VAE-based music clustering pipeline that:
\begin{itemize}
    \item Extracts audio features from 30-second music clips using log-mel spectrograms and MFCCs
    \item Incorporates lyrics embeddings using TF-IDF vectorization for multi-modal analysis
    \item Implements multiple VAE architectures: standard VAE, Convolutional VAE, Beta-VAE ($\beta=4.0$), and Conditional VAE
    \item Compares various clustering algorithms: K-Means, Agglomerative (Ward linkage), DBSCAN, and Gaussian Mixture Models
    \item Performs multi-modal fusion of audio, lyrics, and genre information with configurable weights
    \item Evaluates performance using both unsupervised metrics (Silhouette, Calinski-Harabasz, Davies-Bouldin) and supervised metrics (ARI, NMI, Purity)
    \item Provides comprehensive visualizations including t-SNE plots, PCA projections, cluster distributions, and genre-cluster heatmaps
\end{itemize}

\section{Related Work}

\subsection{Variational Autoencoders}
VAEs, introduced by Kingma and Welling (2014), combine neural networks with variational inference to learn latent representations. The model consists of an encoder $q_\phi(z|x)$ that maps inputs to a latent distribution, and a decoder $p_\theta(x|z)$ that reconstructs inputs from latent samples. The training objective maximizes the Evidence Lower Bound (ELBO):
\begin{equation}
\mathcal{L}(\theta, \phi; x) = \mathbb{E}_{q_\phi(z|x)}[\log p_\theta(x|z)] - D_{KL}(q_\phi(z|x) \| p(z))
\end{equation}

\subsection{Beta-VAE for Disentanglement}
Higgins et al. (2017) introduced Beta-VAE, which adds a hyperparameter $\beta > 1$ to the KL divergence term, encouraging more disentangled latent representations:
\begin{equation}
\mathcal{L}_{\beta} = \mathbb{E}_{q_\phi(z|x)}[\log p_\theta(x|z)] - \beta \cdot D_{KL}(q_\phi(z|x) \| p(z))
\end{equation}

\subsection{Music Representation Learning}
Audio features such as Mel-Frequency Cepstral Coefficients (MFCCs) and mel-spectrograms have been widely used for music analysis. Recent work has explored combining audio with textual features like lyrics for richer multi-modal representations.

\section{Method}

\subsection{Feature Extraction}

\subsubsection{Audio Features}
We extract log-mel spectrogram features from 30-second audio clips:
\begin{itemize}
    \item Sample rate: 22,050 Hz
    \item Number of mel bands: 64
    \item Hop length: 512 samples
    \item Features: concatenation of mean and standard deviation across time (128 dimensions)
\end{itemize}

\subsubsection{Lyrics Features}
For clips with available lyrics, we extract TF-IDF features:
\begin{itemize}
    \item Maximum vocabulary size: 500 words
    \item N-gram range: (1, 2) for unigrams and bigrams
    \item Dimensionality reduction via SVD to 64 dimensions
    \item Additional statistics: word count, unique word ratio, average word length
\end{itemize}

\subsection{VAE Architectures}

\subsubsection{Standard VAE}
Our basic VAE uses fully-connected layers:
\begin{itemize}
    \item Encoder: Input (128) $\rightarrow$ FC (256) $\rightarrow$ ReLU $\rightarrow$ FC (128) $\rightarrow$ $\mu$, $\log\sigma^2$ (32 each)
    \item Decoder: Latent (32) $\rightarrow$ FC (128) $\rightarrow$ ReLU $\rightarrow$ FC (256) $\rightarrow$ ReLU $\rightarrow$ FC (128)
    \item Latent dimension: 32
\end{itemize}

\subsubsection{Convolutional VAE}
For 2D spectrogram inputs, we use convolutional layers:
\begin{itemize}
    \item Encoder: 2D convolutions with stride 2 for downsampling
    \item Decoder: Transposed convolutions for upsampling
    \item Better captures local temporal and frequency patterns
\end{itemize}

\subsubsection{Beta-VAE}
We set $\beta = 4.0$ to encourage disentanglement while maintaining reconstruction quality. Higher beta values force the model to learn more independent latent factors.

\subsubsection{Conditional VAE (CVAE)}
Genre information is incorporated as conditioning:
\begin{itemize}
    \item Genre labels are one-hot encoded
    \item Concatenated with input for encoder and latent code for decoder
    \item Allows genre-aware latent representations
\end{itemize}

\subsection{Clustering Algorithms}

We evaluate four clustering algorithms on the learned latent representations:

\begin{itemize}
    \item \textbf{K-Means}: Partition-based clustering minimizing within-cluster variance
    \item \textbf{Agglomerative Clustering}: Hierarchical clustering with Ward linkage
    \item \textbf{DBSCAN}: Density-based clustering for arbitrary cluster shapes
    \item \textbf{Gaussian Mixture Model}: Probabilistic clustering assuming Gaussian distributions
\end{itemize}

\subsection{Multi-modal Fusion}
For the hard task, we combine audio, lyrics, and genre features using weighted early fusion:
\begin{equation}
f_{fused} = w_a \cdot f_{audio} \oplus w_l \cdot f_{lyrics} \oplus w_g \cdot f_{genre}
\end{equation}
where $\oplus$ denotes concatenation and weights are $(w_a, w_l, w_g) = (0.6, 0.3, 0.1)$.

\section{Experiments}

\subsection{Dataset}
Our dataset is a hybrid language music collection containing both English and Bangla songs:
\begin{itemize}
    \item \textbf{Total clips}: 3,213 audio segments (30 seconds each)
    \item \textbf{Genres}: 82 unique genre labels
    \item \textbf{Feature dimension}: 128 (concatenated mean and std of log-mel spectrograms)
    \item \textbf{Latent dimension}: 32
    \item \textbf{Number of clusters}: 4 (determined via elbow method analysis)
\end{itemize}

\subsubsection{Genre Distribution}
The dataset exhibits a long-tail distribution across genres. The top 20 genres by frequency are:

\begin{table}[h]
\caption{Genre Distribution in Dataset (Top 20)}
\label{tab:genres}
\centering
\small
\begin{tabular}{lc|lc}
\toprule
Genre & Count & Genre & Count \\
\midrule
R\&B/Soul & 593 & Alternative & 62 \\
Metal & 299 & New Wave & 50 \\
Hip Hop & 262 & Reggae & 41 \\
Pop Rock & 247 & Rock & 37 \\
Alternative Indie & 198 & Funk & 36 \\
Country & 187 & Hard Rock & 31 \\
Alternative Rock & 161 & Pop Rap & 28 \\
Pop & 142 & Dance Pop & 25 \\
Synth-pop & 109 & R\&B & 24 \\
Soft Rock & 84 & Country Rock & 24 \\
\bottomrule
\end{tabular}
\end{table}

The remaining 62 genres each have fewer than 25 clips, including Contemporary R\&B (21), Trap (19), Reggae-Pop (19), Country Pop (18), Folk (18), Dance/Electronic (17), Folk Rock (15), UK Rap (13), Nu Metal (13), Punk Rock (11), Post-Punk (11), Regional Mexican (10), Soul Folk (10), and many others.

\subsection{Training Details}
\begin{itemize}
    \item Optimizer: Adam with learning rate 0.001
    \item Epochs: 50
    \item Batch size: 64
    \item Number of clusters: 4 (determined via elbow method analysis)
    \item Random seed: 42 for reproducibility
\end{itemize}

\subsection{Evaluation Metrics}

\subsubsection{Unsupervised Metrics}
\begin{itemize}
    \item \textbf{Silhouette Score}: Measures cluster cohesion and separation. Range: [-1, 1], higher is better.
    \begin{equation}
    s(i) = \frac{b(i) - a(i)}{\max(a(i), b(i))}
    \end{equation}
    
    \item \textbf{Calinski-Harabasz Index}: Ratio of between-cluster to within-cluster variance. Higher is better.
    
    \item \textbf{Davies-Bouldin Index}: Average cluster similarity. Lower is better.
\end{itemize}

\subsubsection{Supervised Metrics (using genre labels)}
\begin{itemize}
    \item \textbf{Adjusted Rand Index (ARI)}: Agreement with ground truth, adjusted for chance. Range: [-1, 1].
    \item \textbf{Normalized Mutual Information (NMI)}: Mutual information normalized to [0, 1].
    \item \textbf{Cluster Purity}: Fraction of dominant class in each cluster.
\end{itemize}

\section{Results}

\subsection{Easy Task: Basic VAE + K-Means}

The Easy Task focuses on implementing a basic VAE for feature extraction and comparing with PCA + K-Means baseline. Table~\ref{tab:easy} presents the results.

\begin{table}[h]
\caption{Easy Task Results: VAE + K-Means vs Baselines}
\label{tab:easy}
\centering
\begin{tabular}{lccccc}
\toprule
Method & Silhouette & CH Index & DB Index & ARI & NMI \\
\midrule
VAE + K-Means & \textbf{0.197} & 704.66 & 1.516 & 0.021 & \textbf{0.103} \\
PCA + K-Means & 0.195 & \textbf{1616.92} & \textbf{1.502} & 0.027 & 0.075 \\
Autoencoder + K-Means & 0.116 & 673.62 & 2.214 & \textbf{0.031} & 0.078 \\
Direct K-Means & 0.183 & 1502.05 & 1.568 & 0.027 & 0.075 \\
\bottomrule
\end{tabular}
\end{table}

\textbf{Analysis:} The basic VAE achieves the highest Silhouette Score (0.197) and NMI (0.103) among all methods. While PCA achieves a higher Calinski-Harabasz Index (1616.92), this metric favors convex clusters and may not capture the complex structure of music data. The VAE's probabilistic latent space provides a richer representation suitable for visualization and generation tasks. The Autoencoder without the variational constraint performs worst, suggesting that the KL divergence regularization is crucial for learning clusterable representations.

\textbf{Visualizations Generated:}
\begin{itemize}
    \item \texttt{easy\_vae\_tsne.png}: t-SNE visualization of VAE latent space colored by cluster
    \item \texttt{easy\_vae\_pca.png}: PCA visualization showing first two principal components
    \item \texttt{easy\_vae\_tsne\_genre.png}: t-SNE colored by genre labels
    \item \texttt{easy\_cluster\_dist.png}: Bar chart of cluster size distribution
    \item \texttt{easy\_genre\_cluster\_heatmap.png}: Heatmap showing genre distribution across clusters
    \item \texttt{easy\_baseline\_comparison.png}: Comparison bar chart of all baseline methods
\end{itemize}

\subsection{Medium Task: Enhanced Architectures}

The Medium Task enhances the VAE with convolutional architecture for spectrograms and includes hybrid audio+lyrics representations. Table~\ref{tab:medium} shows results with multiple clustering algorithms.

\begin{table}[h]
\caption{Medium Task Results: ConvVAE, Lyrics, and Multiple Clustering Algorithms}
\label{tab:medium}
\centering
\begin{tabular}{lcccc}
\toprule
Method & Silhouette & DB Index & ARI & NMI \\
\midrule
ConvVAE + K-Means & \textbf{0.333} & 0.844 & \textbf{0.066} & \textbf{0.209} \\
Lyrics TF-IDF + K-Means & 0.328 & \textbf{0.751} & 0.013 & 0.078 \\
VAE + DBSCAN & 0.249 & 2.262 & 0.002 & 0.012 \\
VAE + K-Means & 0.169 & 1.645 & 0.027 & 0.106 \\
VAE + Agglomerative & 0.138 & 1.927 & 0.011 & 0.084 \\
VAE + GMM & 0.131 & 1.908 & 0.025 & 0.097 \\
VAE + Lyrics Combined & 0.114 & 1.587 & 0.034 & 0.095 \\
\bottomrule
\end{tabular}
\end{table}

\textbf{Analysis of Clustering Algorithms:}
\begin{itemize}
    \item \textbf{ConvVAE + K-Means}: Achieves the best overall performance with Silhouette 0.333 and NMI 0.209. The convolutional architecture better captures local temporal and frequency patterns in spectrograms compared to fully-connected networks.
    
    \item \textbf{Lyrics TF-IDF + K-Means}: Surprisingly competitive with Silhouette 0.328, demonstrating that textual content captures meaningful genre distinctions. The lowest Davies-Bouldin Index (0.751) indicates well-separated clusters.
    
    \item \textbf{DBSCAN}: Higher Silhouette (0.249) than standard K-Means on VAE features, but very low ARI/NMI, suggesting it finds density-based clusters that don't align with genre labels.
    
    \item \textbf{Agglomerative \& GMM}: Lower performance than K-Means, possibly because the latent space structure is better suited to spherical clusters.
    
    \item \textbf{Audio + Lyrics Combined}: Concatenating audio and lyrics features surprisingly decreased performance, suggesting naive fusion is suboptimal.
\end{itemize}

\textbf{Why VAE Representations Perform Better/Worse:}
The VAE's KL divergence constraint encourages a smooth, Gaussian-like latent space, which is well-suited for K-Means but may not be optimal for density-based (DBSCAN) or hierarchical methods. The regularization prevents the encoder from creating arbitrary cluster shapes, which explains why K-Means consistently outperforms other algorithms on VAE features.

\subsection{Hard Task: Advanced VAE Variants}

The Hard Task implements Beta-VAE and CVAE for disentangled representations, and multi-modal fusion combining audio, lyrics, and genre. Table~\ref{tab:hard} presents the results.

\begin{table}[h]
\caption{Hard Task Results: Beta-VAE, CVAE, and Multi-modal Fusion}
\label{tab:hard}
\centering
\begin{tabular}{lcccccc}
\toprule
Method & Silhouette & CH Index & DB Index & ARI & NMI & Purity \\
\midrule
Beta-VAE + K-Means & \textbf{0.366} & \textbf{2375.42} & \textbf{0.974} & 0.023 & 0.076 & \textbf{0.228} \\
CVAE + K-Means & 0.237 & 873.15 & 1.349 & -0.005 & 0.034 & 0.185 \\
Multi-modal Fusion & 0.070 & 374.35 & 2.684 & 0.016 & 0.097 & 0.215 \\
\bottomrule
\end{tabular}
\end{table}

\textbf{Analysis of Advanced Methods:}
\begin{itemize}
    \item \textbf{Beta-VAE ($\beta=4.0$)}: Achieves the best Silhouette Score (0.366) across all experiments. The disentanglement constraint forces the model to learn independent latent factors, resulting in cleaner cluster boundaries. The high Calinski-Harabasz Index (2375.42) confirms well-separated, compact clusters.
    
    \item \textbf{Conditional VAE}: Genre conditioning surprisingly hurts clustering performance (negative ARI). This suggests that conditioning on noisy or overly fine-grained labels (82 genres) may confuse the model rather than help.
    
    \item \textbf{Multi-modal Fusion}: Concatenating audio, lyrics, and genre features with weights $(0.6, 0.3, 0.1)$ yields the lowest Silhouette Score. This indicates that simple early fusion is not effective; more sophisticated fusion strategies (attention-based, late fusion) are needed.
\end{itemize}

\textbf{Cluster Purity Analysis:}
Overall cluster purity is low (22.8\% for Beta-VAE) because we use only 4 clusters for 82 genres. However, examination of individual clusters reveals meaningful groupings:
\begin{itemize}
    \item Clusters tend to group similar genres (e.g., R\&B/Soul often clusters with Hip Hop)
    \item Rock-related genres (Metal, Hard Rock, Alternative Rock) tend to cluster together
    \item Pop variants (Pop, Dance Pop, Synth-pop) show clustering affinity
\end{itemize}

\subsection{Overall Comparison}

Table~\ref{tab:final} provides a comprehensive ranking of all methods across all task levels.

\begin{table}[h]
\caption{Final Comparison: All 14 Methods Ranked by Silhouette Score}
\label{tab:final}
\centering
\small
\begin{tabular}{lcccc}
\toprule
Method & Silhouette & ARI & NMI & Purity \\
\midrule
Beta-VAE + K-Means & \textbf{0.366} & 0.023 & 0.076 & 0.228 \\
ConvVAE + K-Means & 0.333 & \textbf{0.066} & \textbf{0.209} & \textbf{0.380} \\
Lyrics TF-IDF + K-Means & 0.328 & 0.013 & 0.078 & 0.186 \\
VAE + DBSCAN & 0.249 & 0.002 & 0.012 & 0.185 \\
CVAE + K-Means & 0.237 & -0.005 & 0.034 & 0.185 \\
VAE + K-Means (Easy) & 0.197 & 0.021 & 0.103 & 0.200 \\
PCA + K-Means & 0.195 & 0.027 & 0.075 & 0.226 \\
Direct K-Means & 0.183 & 0.027 & 0.075 & 0.227 \\
VAE + K-Means (Medium) & 0.169 & 0.027 & 0.106 & 0.210 \\
VAE + Agglomerative & 0.138 & 0.011 & 0.084 & 0.198 \\
VAE + GMM & 0.131 & 0.025 & 0.097 & 0.201 \\
Autoencoder + K-Means & 0.116 & 0.031 & 0.078 & 0.230 \\
VAE + Lyrics Combined & 0.114 & 0.034 & 0.095 & 0.216 \\
Multi-modal Fusion & 0.070 & 0.016 & 0.097 & 0.215 \\
\bottomrule
\end{tabular}
\end{table}

\textbf{Key Observations:}
\begin{enumerate}
    \item \textbf{Best by Silhouette Score}: Beta-VAE + K-Means (0.366)
    \item \textbf{Best by ARI}: ConvVAE + K-Means (0.066)
    \item \textbf{Best by NMI}: ConvVAE + K-Means (0.209)
    \item \textbf{Best by Purity}: ConvVAE + K-Means (0.380)
\end{enumerate}

The visualizations confirm these quantitative findings. Figure analysis shows that Beta-VAE produces the most well-separated clusters in t-SNE space, while ConvVAE clusters align better with ground-truth genre labels.

\section{Discussion}

\subsection{Key Findings}
\begin{enumerate}
    \item \textbf{Beta-VAE achieves best cluster separation}: The disentanglement constraint ($\beta=4$) encourages learning independent latent factors, resulting in cleaner cluster boundaries. The Silhouette Score improvement from 0.197 (basic VAE) to 0.366 (Beta-VAE) represents an 86\% improvement.
    
    \item \textbf{Convolutional architecture excels at supervised metrics}: ConvVAE captures local temporal and frequency patterns in spectrograms, leading to the highest NMI (0.209) and Purity (0.380). This suggests that 2D convolutional processing better preserves genre-discriminative information.
    
    \item \textbf{Lyrics provide complementary but challenging information}: TF-IDF features from lyrics achieve competitive Silhouette Scores (0.328), suggesting textual content captures meaningful genre distinctions. However, combining audio and lyrics features naively (concatenation) degrades performance, indicating the need for more sophisticated fusion strategies.
    
    \item \textbf{K-Means consistently outperforms other clustering algorithms}: On VAE latent spaces, K-Means achieves better results than Agglomerative, DBSCAN, and GMM. This is expected because the KL divergence regularization encourages Gaussian-like distributions, which are well-suited for centroid-based clustering.
    
    \item \textbf{Multi-modal fusion is challenging}: Simple weighted concatenation of audio, lyrics, and genre features did not outperform single-modality approaches. The modalities have different scales, semantics, and noise characteristics that require careful alignment.
    
    \item \textbf{Low ARI/NMI reflects genre complexity}: With 82 genres and only 4 clusters, exact alignment with labels is not expected. The metrics reflect that clusters capture broader musical similarities (e.g., "Rock-like" vs "R\&B-like") rather than fine-grained genre distinctions.
    
    \item \textbf{Baseline PCA is competitive but limited}: PCA + K-Means achieves similar Silhouette (0.195) to basic VAE (0.197), but VAE provides a generative model capable of sampling and interpolation, which PCA cannot offer.
\end{enumerate}

\subsection{Comparison with Requirements}

Table~\ref{tab:requirements} summarizes how our implementation addresses each task requirement:

\begin{table}[h]
\caption{Implementation Coverage of Task Requirements}
\label{tab:requirements}
\centering
\small
\begin{tabular}{p{6cm}c}
\toprule
Requirement & Status \\
\midrule
\textbf{Easy Task} & \\
Basic VAE for feature extraction & \checkmark \\
Hybrid language dataset (English + Bangla) & \checkmark \\
K-Means clustering on latent features & \checkmark \\
t-SNE visualization & \checkmark \\
PCA visualization & \checkmark \\
Silhouette Score evaluation & \checkmark \\
Calinski-Harabasz Index evaluation & \checkmark \\
Baseline comparison (PCA + K-Means) & \checkmark \\
\midrule
\textbf{Medium Task} & \\
Convolutional VAE for spectrograms & \checkmark \\
Hybrid feature: audio + lyrics embeddings & \checkmark \\
K-Means clustering & \checkmark \\
Agglomerative Clustering & \checkmark \\
DBSCAN clustering & \checkmark \\
Davies-Bouldin Index evaluation & \checkmark \\
Adjusted Rand Index (ARI) & \checkmark \\
Analysis of VAE vs baseline performance & \checkmark \\
\midrule
\textbf{Hard Task} & \\
Beta-VAE for disentanglement & \checkmark \\
Conditional VAE (CVAE) & \checkmark \\
Multi-modal: audio + lyrics + genre & \checkmark \\
Normalized Mutual Information (NMI) & \checkmark \\
Cluster Purity & \checkmark \\
Latent space visualizations & \checkmark \\
Cluster distribution over genres & \checkmark \\
Comparison: VAE vs PCA vs Autoencoder vs Direct & \checkmark \\
\bottomrule
\end{tabular}
\end{table}

\subsection{Limitations}
\begin{itemize}
    \item The number of clusters (k=4) may not capture the full diversity of 82 genres. Hierarchical or variable-k clustering could be explored.
    \item Clip-level lyrics may be incomplete or noisy, especially for songs with instrumental sections.
    \item Training was performed without extensive hyperparameter tuning due to computational constraints.
    \item The dataset is relatively small (3,213 clips) compared to large-scale music datasets like Million Song Dataset.
    \item The hybrid language aspect (English + Bangla) could be further explored with language-specific models.
\end{itemize}

\subsection{Future Work}
\begin{itemize}
    \item Explore hierarchical clustering to capture genre hierarchies and sub-genres
    \item Implement attention-based multi-modal fusion for better cross-modal alignment
    \item Use pre-trained audio embeddings (e.g., from music transformers like Jukebox or MusicLM)
    \item Investigate semi-supervised approaches with partial labels to improve genre alignment
    \item Apply UMAP as an alternative to t-SNE for visualization
    \item Explore contrastive learning objectives for better cluster separation
\end{itemize}

\section{Conclusion}

We presented a comprehensive study of VAE-based unsupervised clustering for hybrid language music data, implementing all three task levels (Easy, Medium, Hard) with 14 different method combinations. Our experiments on 3,213 audio clips spanning 82 genres demonstrate several key findings:

\begin{enumerate}
    \item \textbf{Beta-VAE} with K-Means clustering achieves the best unsupervised performance with a Silhouette Score of 0.366, an 86\% improvement over the basic VAE baseline (0.197).
    
    \item \textbf{Convolutional VAE} achieves the best supervised metrics (NMI=0.209, Purity=0.380), demonstrating that 2D processing of spectrograms preserves genre-discriminative information.
    
    \item \textbf{Lyrics embeddings} provide competitive clustering performance (Silhouette=0.328), confirming that textual content captures meaningful musical distinctions.
    
    \item \textbf{K-Means} consistently outperforms Agglomerative, DBSCAN, and GMM on VAE latent spaces due to the Gaussian-like structure induced by KL regularization.
    
    \item \textbf{Multi-modal fusion} requires sophisticated strategies; simple concatenation degrades performance.
\end{enumerate}

While exact alignment with fine-grained genre labels remains challenging (82 genres into 4 clusters), the learned clusters capture meaningful musical similarities that can support applications like music recommendation, playlist generation, and music discovery. The VAE framework additionally enables generative capabilities like latent space interpolation and sample generation, which linear methods like PCA cannot offer.

All code, data, and visualizations are available for reproducibility. The project successfully addresses all requirements across Easy, Medium, and Hard task levels.

\section*{Acknowledgments}
This project was completed as part of CSE425 Neural Networks course at BRAC University.

\bibliographystyle{plain}
\begin{thebibliography}{9}

\bibitem{kingma2014}
Kingma, D.P. and Welling, M. (2014).
Auto-Encoding Variational Bayes.
\textit{International Conference on Learning Representations (ICLR)}.

\bibitem{higgins2017}
Higgins, I., et al. (2017).
beta-VAE: Learning Basic Visual Concepts with a Constrained Variational Framework.
\textit{International Conference on Learning Representations (ICLR)}.

\bibitem{sohn2015}
Sohn, K., Lee, H., and Yan, X. (2015).
Learning Structured Output Representation using Deep Conditional Generative Models.
\textit{Advances in Neural Information Processing Systems (NeurIPS)}.

\bibitem{mcfee2015}
McFee, B., et al. (2015).
librosa: Audio and Music Signal Analysis in Python.
\textit{Proceedings of the 14th Python in Science Conference}.

\bibitem{sklearn}
Pedregosa, F., et al. (2011).
Scikit-learn: Machine Learning in Python.
\textit{Journal of Machine Learning Research}, 12:2825-2830.

\end{thebibliography}

\appendix
\section{Appendix A: Generated Visualizations}

The following visualizations were generated during the experiments and are available in the \texttt{output5/} directory:

\subsection{Easy Task Visualizations}
\begin{enumerate}
    \item \texttt{easy\_vae\_tsne.png} - t-SNE visualization of VAE latent space colored by cluster assignment
    \item \texttt{easy\_vae\_pca.png} - PCA visualization showing first two principal components with explained variance
    \item \texttt{easy\_vae\_tsne\_genre.png} - t-SNE colored by ground-truth genre labels
    \item \texttt{easy\_cluster\_dist.png} - Bar chart showing cluster size distribution
    \item \texttt{easy\_genre\_cluster\_heatmap.png} - Heatmap showing genre distribution across clusters
    \item \texttt{easy\_baseline\_comparison.png} - Bar chart comparing all baseline methods
\end{enumerate}

\subsection{Medium Task Visualizations}
\begin{enumerate}
    \item \texttt{medium\_kmeans\_tsne.png} - t-SNE of VAE + K-Means clustering
    \item \texttt{medium\_agglomerative\_tsne.png} - t-SNE of VAE + Agglomerative clustering
    \item \texttt{medium\_dbscan\_tsne.png} - t-SNE of VAE + DBSCAN clustering
    \item \texttt{medium\_gmm\_tsne.png} - t-SNE of VAE + GMM clustering
    \item \texttt{medium\_combined\_tsne.png} - t-SNE of Audio + Lyrics combined features
    \item \texttt{medium\_clustering\_comparison.png} - Comparison of all clustering algorithms
\end{enumerate}

\subsection{Hard Task Visualizations}
\begin{enumerate}
    \item \texttt{hard\_beta\_vae\_tsne.png} - Beta-VAE latent space visualization by cluster
    \item \texttt{hard\_beta\_vae\_genre.png} - Beta-VAE latent space colored by genre
    \item \texttt{hard\_beta\_vae\_heatmap.png} - Genre-cluster distribution heatmap for Beta-VAE
    \item \texttt{hard\_cvae\_tsne.png} - CVAE latent space by genre conditioning
    \item \texttt{hard\_multimodal\_tsne.png} - Multi-modal fusion t-SNE visualization
    \item \texttt{hard\_multimodal\_heatmap.png} - Genre-cluster distribution for multi-modal fusion
    \item \texttt{hard\_methods\_comparison.png} - Comparison of all hard task methods
    \item \texttt{final\_comparison.png} - Bar chart comparing all 14 methods across all tasks
\end{enumerate}

\section{Appendix B: Hyperparameters}

\begin{table}[h]
\centering
\begin{tabular}{ll}
\toprule
Parameter & Value \\
\midrule
\multicolumn{2}{l}{\textbf{Audio Feature Extraction}} \\
Sample Rate & 22,050 Hz \\
Mel Bands & 64 \\
Hop Length & 512 samples \\
Clip Duration & 30 seconds \\
\midrule
\multicolumn{2}{l}{\textbf{VAE Architecture}} \\
Input Dimension & 128 \\
Hidden Dimension & 256 \\
Latent Dimension & 32 \\
Activation & ReLU \\
\midrule
\multicolumn{2}{l}{\textbf{Training}} \\
Optimizer & Adam \\
Learning Rate & 0.001 \\
Batch Size & 64 \\
Epochs & 50 \\
Random Seed & 42 \\
\midrule
\multicolumn{2}{l}{\textbf{Advanced VAE}} \\
Beta (Beta-VAE) & 4.0 \\
CVAE Genre Classes & 82 \\
\midrule
\multicolumn{2}{l}{\textbf{Lyrics Features}} \\
TF-IDF Max Features & 500 \\
N-gram Range & (1, 2) \\
SVD Components & 64 \\
\midrule
\multicolumn{2}{l}{\textbf{Multi-modal Fusion}} \\
Audio Weight & 0.6 \\
Lyrics Weight & 0.3 \\
Genre Weight & 0.1 \\
\midrule
\multicolumn{2}{l}{\textbf{Clustering}} \\
Number of Clusters & 4 \\
K-Means n\_init & 10 \\
DBSCAN eps & 0.5 \\
DBSCAN min\_samples & 5 \\
\bottomrule
\end{tabular}
\caption{Complete hyperparameter configuration}
\end{table}

\section{Appendix C: Cluster Listings}

Detailed cluster assignments are saved in the following files:
\begin{itemize}
    \item \texttt{clusters\_easy\_vae\_kmeans.txt} - VAE + K-Means cluster assignments
    \item \texttt{clusters\_easy\_pca.txt} - PCA + K-Means cluster assignments
    \item \texttt{clusters\_easy\_autoencoder.txt} - Autoencoder + K-Means assignments
    \item \texttt{clusters\_easy\_direct.txt} - Direct K-Means assignments
    \item \texttt{clusters\_medium\_*.txt} - Medium task cluster assignments
    \item \texttt{clusters\_hard\_*.txt} - Hard task cluster assignments
    \item \texttt{cluster\_listing.txt} - Master cluster listing with all clips
    \item \texttt{all\_metrics.json} - JSON file with all computed metrics
\end{itemize}

\end{document}
